\chapter{Verb Position and the Sentence Bracket}

German word order becomes much easier when the \textbf{finite verb} is found
first. The finite verb is the form that agrees with the subject and carries
the tense: \textit{lernt}, \textit{hat}, \textit{muss}, or \textit{wird}.
Other verb parts often stand at the end of the clause and form a
\textit{Satzklammer} (sentence bracket) with it.

\section{Main clauses: the finite verb is in position two}

In a statement, the finite verb is the second \emph{constituent}, not
necessarily the second word. Any one complete element can occupy position one.

\begin{center}
\small
\rowcolors{2}{referenceBlueBg}{white}
\begin{tabularx}{\textwidth}{@{}>{\bfseries}l l l X@{}}
\toprule
Position 1 & Finite verb & Middle field & End field \\
\midrule
Ich & \textbf{lerne} & heute Abend & Deutsch. \\
Heute Abend & \textbf{lerne} & ich & Deutsch. \\
Nach dem Abendessen & \textbf{lerne} & ich zu Hause & Deutsch. \\
\bottomrule
\end{tabularx}
\end{center}

Only one element comes before the finite verb. If something other than the
subject occupies position one, the subject normally follows the verb:
\textit{Morgen \textbf{fahre} ich nach Berlin}, not
\textit{Morgen ich fahre nach Berlin}.

\section{The Satzklammer}

With separable verbs, compound tenses, modal verbs, the future, and the
passive, the finite verb opens the bracket and the remaining verb material
closes it.

\begin{center}
\small
\rowcolors{2}{referenceGreenBg}{white}
\begin{tabularx}{\textwidth}{@{}>{\bfseries}l l X l@{}}
\toprule
Construction & Left bracket & Middle field & Right bracket \\
\midrule
Separable verb & Ich \textbf{rufe} & dich heute Abend & \textbf{an}. \\
Perfekt & Ich \textbf{habe} & dich gestern & \textbf{angerufen}. \\
Modal verb & Ich \textbf{muss} & dich heute Abend & \textbf{anrufen}. \\
Futur I & Ich \textbf{werde} & dich morgen & \textbf{anrufen}. \\
Passive & Das Formular \textbf{wird} & heute & \textbf{ausgefüllt}. \\
\bottomrule
\end{tabularx}
\end{center}

Material such as objects, adverbs, and negation usually stands inside the
bracket. The right bracket should not be forgotten merely because the middle
field is long.

\section{Placement of \textit{gern(e)}, \textit{nicht}, and \textit{nichts}}

\subsection{\textit{gern} and \textit{gerne}}

\textit{Gern} and \textit{gerne} have the same meaning and are both standard.
They show that somebody likes doing an activity. In the neutral word order,
\textit{gern(e)} stands in the middle field, usually after the finite verb and
pronouns but before the final verb element.

\begin{center}
\small
\rowcolors{2}{referenceBlueBg}{white}
\begin{tabularx}{\textwidth}{@{}>{\bfseries}l X l@{}}
\toprule
Construction & Example & Meaning \\
\midrule
Simple verb & Ich lese \textbf{gern} Krimis. & I like reading crime novels. \\
Modal verb & Ich möchte \textbf{gerne} mitkommen. & I would like to come along. \\
Separable verb & Ich komme \textbf{gern} mit. & I am happy to come along. \\
Perfekt & Ich habe dort \textbf{gern} gearbeitet. & I liked working there. \\
Short answer & Kommst du mit? -- Ja, \textbf{gern}! & Yes, gladly. \\
\bottomrule
\end{tabularx}
\end{center}

The comparison forms are \textit{lieber} and \textit{am liebsten}:
\textit{Ich fahre lieber mit dem Zug};
\textit{Am liebsten bleibe ich zu Hause}.

\subsection{Choosing \textit{nicht}, \textit{nichts}, or \textit{kein}}

These forms do different grammatical jobs.

\begin{center}
\small
\rowcolors{2}{referenceGreenBg}{white}
\begin{tabularx}{\textwidth}{@{}>{\bfseries}l l X@{}}
\toprule
Form & Meaning and function & Example \\
\midrule
nicht & not; negates an action, quality, or particular element
      & \textit{Das ist nicht teuer.} \\
nichts & nothing; stands in place of a noun or object
       & \textit{Ich sehe nichts.} \\
kein & no/not a; accompanies a noun and declines like \textit{ein}
     & \textit{Ich habe keine Zeit.} \\
\bottomrule
\end{tabularx}
\end{center}

To negate one particular element, put \textit{nicht} directly before it:
\textit{nicht heute}, \textit{nicht schnell}, \textit{nicht mit dem Bus},
\textit{nicht nach Berlin}. This position often creates a contrast:
\textit{Ich komme nicht heute, sondern morgen}.

To negate the action or clause as a whole, \textit{nicht} normally stands as
late as possible. It comes before the closing part of a sentence bracket.

\begin{center}
\small
\rowcolors{2}{referenceYellowBg}{white}
\begin{tabularx}{\textwidth}{@{}>{\bfseries}l X@{}}
\toprule
Construction & Example \\
\midrule
Simple verb & \textit{Ich arbeite heute \textbf{nicht}.} \\
Definite object & \textit{Ich kaufe das Auto \textbf{nicht}.} \\
Separable verb & \textit{Ich rufe ihn heute \textbf{nicht} an.} \\
Modal verb & \textit{Ich kann heute \textbf{nicht} kommen.} \\
Perfekt & \textit{Ich habe den Film \textbf{nicht} gesehen.} \\
\bottomrule
\end{tabularx}
\end{center}

\textit{Nichts} occupies an ordinary object position:
\textit{Ich weiß \textbf{nichts} darüber};
\textit{Ich habe \textbf{nichts} verstanden}. Use \textit{kein} with an
indefinite or article-free noun: \textit{Ich trinke \textbf{keinen} Kaffee};
\textit{Sie hat \textbf{keine} Kinder}.

The position changes the meaning:
\textit{Ich koche \textbf{nicht gern}} means \enquote{I do not like cooking}, while
\textit{Ich koche heute \textbf{nicht}} means \enquote{I am not cooking today}.
Other useful combinations are \textit{noch nicht} (not yet),
\textit{nicht mehr} (not any more), and \textit{gar nicht} (not at all).

\section{Subordinate clauses: the finite verb moves to the end}

After a subordinating conjunction such as \textit{weil, dass, obwohl}, or
\textit{wenn}, the finite verb normally stands at the end of its clause. A
separable verb is written as one word there.

\begin{center}
\small
\rowcolors{2}{referenceYellowBg}{white}
\begin{tabularx}{\textwidth}{@{}>{\bfseries}l X@{}}
\toprule
Construction & Subordinate clause \\
\midrule
Simple verb & \ldots{}, weil er heute lange \textbf{arbeitet}. \\
Separable verb & \ldots{}, weil er mich später \textbf{anruft}. \\
Perfekt & \ldots{}, weil er mich gestern \textbf{angerufen hat}. \\
Modal verb & \ldots{}, weil er morgen lange \textbf{arbeiten muss}. \\
Passive & \ldots{}, weil das Formular heute \textbf{ausgefüllt wird}. \\
\bottomrule
\end{tabularx}
\end{center}

Relative clauses follow the same verb-final pattern:
\textit{Das ist die Frau, die mir gestern \textbf{geholfen hat}.}

When the subordinate clause comes first, the complete clause occupies
position one. The finite verb of the following main clause therefore comes
immediately after the comma:

\begin{center}
\textit{Wenn ich Zeit \textbf{habe}, \textbf{rufe} ich dich an.}
\end{center}

\section{Questions and commands}

Yes/no questions and commands normally begin with the finite verb. A question
word, however, occupies position one in an information question.

\begin{center}
\small
\rowcolors{2}{referenceRedBg}{white}
\begin{tabularx}{\textwidth}{@{}>{\bfseries}l l X@{}}
\toprule
Type & Pattern & Example \\
\midrule
Yes/no question & finite verb + subject & \textbf{Kommst} du morgen? \\
Information question & question word + finite verb & Wann \textbf{kommst} du? \\
Command & imperative + remaining information & \textbf{Ruf} mich morgen \textbf{an}! \\
Formal command & verb + \textit{Sie} & \textbf{Rufen Sie} mich morgen \textbf{an}! \\
\bottomrule
\end{tabularx}
\end{center}

\begin{tcolorbox}[
  colback=referenceBlueBg,
  colframe=genderMasculine,
  title={Quick verb-position check},
  fonttitle=\bfseries
]
\begin{enumerate}[leftmargin=*,nosep]
  \item Find the finite verb.
  \item Decide whether the clause is a main clause, question, command, or
        subordinate clause.
  \item Check whether another verb part must close the sentence bracket.
  \item In a main-clause statement, make sure exactly one element precedes the
        finite verb.
\end{enumerate}
\end{tcolorbox}
